% Chapter 4: Analysis

\section{Introduction}

This chapter presents the analysis phase of the Lesaka AI project, including functional and non-functional requirements, use cases, and system constraints. The analysis establishes the foundation for system design and implementation.

\section{Functional Requirements}

\subsection{User Management}

\begin{itemize}
    \item \textbf{FR-001:} The system shall allow farmers to register with their personal details
    \item \textbf{FR-002:} The system shall generate anonymized owner IDs for privacy protection
    \item \textbf{FR-003:} The system shall support multiple user roles (admin, farmer, broker)
    \item \textbf{FR-004:} The system shall implement role-based access control
\end{itemize}

\subsection{Cattle Management}

\begin{itemize}
    \item \textbf{FR-005:} The system shall allow registration of cattle with unique identifiers
    \item \textbf{FR-006:} The system shall store cattle breed, age, and health information
    \item \textbf{FR-007:} The system shall allow farmers to view their registered cattle
    \item \textbf{FR-008:} The system shall support cattle search and filtering
\end{itemize}

\subsection{Telemetry Management}

\begin{itemize}
    \item \textbf{FR-009:} The system shall accept telemetry data from sensors
    \item \textbf{FR-010:} The system shall validate temperature readings (30.0-45.0°C)
    \item \textbf{FR-011:} The system shall store validated telemetry data
    \item \textbf{FR-012:} The system shall provide telemetry history for individual cattle
\end{itemize}

\subsection{Agent Orchestration}

\begin{itemize}
    \item \textbf{FR-013:} The system shall route telemetry to appropriate agents
    \item \textbf{FR-014:} The system shall implement Agent Molemo for fever detection
    \item \textbf{FR-015:} The system shall implement Agent Loapi for environmental analysis
    \item \textbf{FR-016:} The system shall implement Agent Thekiso for market assessment
\end{itemize}

\section{Non-Functional Requirements}

\subsection{Performance}

\begin{itemize}
    \item \textbf{NFR-001:} The system shall respond to API requests within 200ms (p95)
    \item \textbf{NFR-002:} The system shall support 1000 concurrent users
    \item \textbf{NFR-003:} The database shall handle 10,000 telemetry records per day
\end{itemize}

\subsection{Security}

\begin{itemize}
    \item \textbf{NFR-004:} The system shall encrypt sensitive data at rest
    \item \textbf{NFR-005:} The system shall implement secure authentication
    \item \textbf{NFR-006:} The system shall log all data access attempts
\end{itemize}

\subsection{Reliability}

\begin{itemize}
    \item \textbf{NFR-007:} The system shall have 99.9\% uptime during testing
    \item \textbf{NFR-008:} The system shall implement self-healing mechanisms
    \item \textbf{NFR-009:} The system shall recover from failures within 30 seconds
\end{itemize}

\subsection{Usability}

\begin{itemize}
    \item \textbf{NFR-010:} The system shall provide an intuitive web interface
    \item \textbf{NFR-011:} The system shall be accessible on mobile devices
    \item \textbf{NFR-012:} The system shall provide clear error messages
\end{itemize}

\section{Business Rules}

\begin{itemize}
    \item \textbf{BR-001:} Only registered farmers can register cattle
    \item \textbf{BR-002:} Temperature readings outside 30.0-45.0°C shall be rejected
    \item \textbf{BR-003:} Broker users cannot access GPS coordinates
    \item \textbf{BR-004:} District must be one of: Orapa, Serowe, Maun, Ghanzi, Francistown, Gaborone, Lobatse
    \item \textbf{BR-005:} Each cattle ID must be unique across the system
\end{itemize}

\section{Assumptions}

\begin{itemize}
    \item Users have basic computer literacy
    \item Internet connectivity is available (though may be intermittent)
    \item Sensor devices provide data in the expected format
    \item Farmers will register with accurate information
\end{itemize}

\section{Constraints}

\subsection{Technical Constraints}

\begin{itemize}
    \item System must run on standard hardware (no specialized equipment)
    \item Database must be SQLite for edge deployment
    \item Frontend must be React-based
    \item Backend must be Python-based
\end{itemize}

\subsection{Time Constraints}

\begin{itemize}
    \item Development must be completed within academic semester timeline
    \item Testing must be completed before submission deadline
    \item Documentation must be finalized before final submission
\end{itemize}

\subsection{Resource Constraints}

\begin{itemize}
    \item Single developer context
    \item Limited access to real-world testing environments
    \item Budget constraints for external services
\end{itemize}

\section{Use Cases}

\subsection{Use Case 1: Farmer Registration}

\begin{itemize}
    \item \textbf{Actor:} Farmer
    \item \textbf{Precondition:} User is not registered
    \item \textbf{Main Flow:}
    \begin{enumerate}
        \item User navigates to registration page
        \item User enters personal details and district
        \item System validates district
        \item System generates anonymized owner ID
        \item System stores farmer information
        \item System confirms successful registration
    \end{enumerate}
    \item \textbf{Postcondition:} User is registered with owner ID
\end{itemize}

\subsection{Use Case 2: Telemetry Ingestion}

\begin{itemize}
    \item \textbf{Actor:} Sensor System
    \item \textbf{Precondition:} Cattle is registered
    \item \textbf{Main Flow:}
    \begin{enumerate}
        \item Sensor transmits telemetry data
        \item System receives data
        \item System validates temperature range
        \item System stores valid data
        \item System routes to appropriate agent
        \item System returns processing result
    \end{enumerate}
    \item \textbf{Postcondition:} Telemetry is stored and processed
\end{itemize}

\subsection{Use Case 3: Cattle Health Monitoring}

\begin{itemize}
    \item \textbf{Actor:} Farmer
    \item \textbf{Precondition:} User is logged in
    \item \textbf{Main Flow:}
    \begin{enumerate}
        \item User views cattle list
        \item User selects specific cattle
        \item System displays health status
        \item System shows telemetry history
        \item System provides agent recommendations
    \end{enumerate}
    \item \textbf{Postcondition:} User has current health information
\end{itemize}

\section{Context Diagram}

The system context identifies external entities and their interactions with the Lesaka AI system:

\begin{itemize}
    \item \textbf{Farmers:} Register cattle, view health status, receive alerts
    \item \textbf{Veterinarians:} Access health data, provide diagnoses
    \item \textbf{Brokers:} View market information (restricted access)
    \item \textbf{Sensors:} Transmit telemetry data
    \item \textbf{Administrators:} Manage system configuration and users
\end{itemize}

\section{Conclusion}

The analysis phase has established clear functional and non-functional requirements, business rules, and use cases for the Lesaka AI system. These requirements align with the module specifications for INFS 401/402 and COMP 401/402, ensuring the system meets academic assessment criteria while addressing real-world agricultural needs in Botswana.
